% !TEX TS-program = xelatex
% !TEX encoding = UTF-8 Unicode
% !Mode:: "TeX:UTF-8"

\documentclass{resume}
\usepackage{zh_CN-Adobefonts_external} % Simplified Chinese Support using external fonts (./fonts/zh_CN-Adobe/)
%\usepackage{zh_CN-Adobefonts_internal} % Simplified Chinese Support using system fonts
% \usepackage{linespacing_fix} % disable extra space before next section
\usepackage{cite}
\usepackage{hyperref} 

\begin{document}
\pagenumbering{gobble} % suppress displaying page number

\name{孙昊哲}

% {E-mail}{mobilephone}{homepage}
% be careful of _ in emaill address
\contactInfo{(+86) 15212796025 }{sunhaozhe347@163.com}{GitHub @expecto347 }{}
% {E-mail}{mobilephone}
% keep the last empty braces!
%\contactInfo{xxx@yuanbin.me}{(+86) 131-221-87xxx}{}

% \section{\faGraduationCap\ 教育背景}
\section{教育背景}
\datedsubsection{\textbf{中国科学技术大学},计算机科学与技术,\textit{工学学士}}{2020.9 - 2024.6}
\datedsubsection{\textbf{查尔姆斯理工大学},计算机科学,\textit{在读硕士研究生}}{2024.9 - 2026.6}

% \section{\faCogs\ IT 技能}
\vspace{1em}
\section{技术能力}
% increase linespacing [parsep=0.5ex]
\begin{itemize}[parsep=0.2ex]
  \item \textbf{编程语言}: Python, C, C++, Go, Java
  \item \textbf{深度学习与机器学习框架}: Pytorch/NumPy/Sklearn/LightGBM/Polar
  \item \textbf{工程工具}: Linux/Docker/Git/MySQL/influxDB/Grafana
  \item \textbf{语言能力}: 英语（托福97），普通话
\end{itemize}

% \end{itemize}
\vspace{1em}
\section{项目经历}
\datedsubsection{\textbf{Deep-EC2：基于深度学习的非短视主动特征获取框架 | 硕士毕业论文}}{2026.01 - 2026.6}
\begin{itemize}
  \item \textbf{双头架构设计}：设计并实现支持动态变长掩码（Mask）输入的双头多层感知机（Dual-Head MLP），将等价类边缘切割（EC2）的非贪心长效决策能力与判别式深度学习的高效性成功融合。
  \item \textbf{算法性能突破}：运用蒙特卡洛重参数化采样技术，通过在特征空间注入噪声评估类别后验变化，成功绕过传统 EC2 依赖全局生成式模型及 $O(|y|^2)$ 级别的计算瓶颈。
  \item \textbf{模型置信度校准}：深度定制带有方差惩罚项的高斯负对数似然损失函数（Gaussian NLL Loss），有效遏制了模型在处理高缺失率早期数据时易产生的“过度自信”问题。
  \item \textbf{工程架构落地}：基于 PyTorch Lightning 搭建工业级两阶段训练 Pipeline（随机掩码预训练 + 端到端强化决策），结合动态衰减的 $\epsilon$-greedy 策略，完美平衡特征空间中的探索与利用。
\end{itemize}

% \datedsubsection{\textbf{基于序列蒙特卡洛（SMC）的动态推荐模型优化 | 美国犹他大学}, 科研助理 (导师: \href{https://www.cs.utah.edu/~zhe/}{Shandian Zhe})}{2023.08-2023.12}
% \begin{itemize}
%   \item \textbf{动态兴趣演化建模}：针对推荐场景中的用户兴趣飘移（Concept Drift）问题，将用户行为建模为非线性状态空间过程，运用序列蒙特卡洛（SMC）算法实现对用户潜在兴趣特征的实时追踪与动态表征。
%   \item \textbf{概率推理与不确定性评估}：独立实现基于粒子滤波（Particle Filtering）的兴趣估计框架，通过维护高维粒子的后验分布，显式捕捉推荐过程中的不确定性，为后续排序阶段的探索与利用（E\&E）平衡提供决策支撑。
%   \item \textbf{模型性能突破}：通过优化粒子重采样策略，在保证计算效率的前提下，有效解决了长尾流量下的兴趣捕获难题，实验结果表明，该方案在召回精度与冷启动物品推荐指标上显著优于传统协同过滤（CF）算法。
% \end{itemize}

\datedsubsection{\textbf{交通标志图像特征提取与自动化分类评估系统 | 新加坡国立大学} (导师: \href{https://tsim17.wixsite.com/terencesim}{Terence Sim})}{2023.07-2023.08}
\begin{itemize}
  \item \textbf{负责交通标志分类系统的核心特征工程与模型训练。}对交通标志数据集进行降噪与归一化等预处理，并提取具有高度鲁棒性的方向梯度直方图（HOG）特征用于后续图像分类任务
  \item 深入评估与对比支持向量机（SVM）、随机森林（Random Forest）及高斯朴素贝叶斯等不同分类器在交通标志识别中的表现，量化分析其对模型准确率与计算时间复杂度的影响
  \item 通过超参数调优与特征筛选策略，在兼顾推理速度的前提下实现系统性能优化，最终基于 SVM 与 HOG 特征组合使识别准确率达到 \textbf{96.67\%}
\end{itemize}

\datedsubsection{\textbf{基于 RFID 信号的高精度定位与轨迹优化系统 | 中国科学技术大学}, 科研助理 (导师: \href{http://staff.ustc.edu.cn/~weigong/}{Wei Gong})}{2022.06-2023.03}
\begin{itemize}
  \item \textbf{参与基于 RFID 的高精度定位系统设计与底层算法研发。}运用马尔可夫决策过程（\textbf{MDP}）与扩展卡尔曼滤波（\textbf{EKF}）对目标物体运动轨迹进行精细化建模，显著提升了定位系统的追踪精度与整体效率
  \item 针对系统初始化阶段进行算法优化，通过改进候选轨迹（Candidate Trajectories）的筛选策略，大幅降低了计算开销并提升了程序的执行效率
  \item 设计并执行系列仿真实验，成功验证了似然函数（Likelihood Function）在特定速度区间及多种复杂运动轨迹下的数学正确性与模型稳定性
\end{itemize}
% \begin{onehalfspacing}
% \end{onehalfspacing}

\vspace{1em}
\section{获奖情况}
% increase linespacing [parsep=0.5ex]
\begin{itemize}[parsep=0.2ex]

\item [2026] Adlerbert 奖学金 (Adlerbert Scholarship)
\item [2024] 瑞典政府 IPOET((International Programme Office for Education and Training) 奖学金
\item [2023] 新加坡国立大学 Visual Computing Summer Workshop一等奖
\item [2020] 中国科学技术大学优秀新生奖学金

\end{itemize}

% \section{\faHeartO\ 项目/作品摘要}
% \section{项目/作品摘要}
% \datedline{\textit{An Integrated Version of Security Monitor Vis System}, https://hijiangtao.github.io/ss-vis-component/ }{}
% \datedline{\textit{Dark-Tech}, https://github.com/hijiangtao/dark-tech/ }{}
% \datedline{\textit{融合社交网络数据挖掘的电视节目可视化分析系统}, https://hijiangtao.github.io/variety-show-hot-spot-vis/}{}
% \datedline{\textit{LeetCodeOJ Solutions}, https://github.com/hijiangtao/LeetCodeOJ}{}
% \datedline{\textit{Info-Vis}, https://github.com/ISCAS-VIS/infovis-ucas}{}

%% Reference
%\newpage
%\bibliographystyle{IEEETran}
%\bibliography{mycite}
\end{document}
